\documentclass[12pt,a4paper]{article}

\usepackage[margin=0.9in]{geometry}
\usepackage{tikz}
\usepackage{booktabs}
\usepackage{tabularx}
\usepackage{colortbl}
\usepackage{multirow}
\usepackage{hyperref}
\usepackage{xcolor}
\usepackage{titlesec}
\usepackage{enumitem}
\usepackage{fancyhdr}
\usepackage{float}
\usepackage{tcolorbox}
\usepackage{listings}
\usepackage{amsmath}
\usepackage{amssymb}

\usetikzlibrary{shapes.geometric, arrows.meta, positioning}
\tcbuselibrary{skins,breakable}

% Colors
\definecolor{primary}{RGB}{0,82,147}
\definecolor{secondary}{RGB}{52,152,219}
\definecolor{accent}{RGB}{39,174,96}
\definecolor{highlight}{RGB}{230,126,34}
\definecolor{danger}{RGB}{231,76,60}
\definecolor{light}{RGB}{236,240,241}
\definecolor{codebg}{RGB}{245,245,245}

% Section formatting
\titleformat{\section}{\Large\bfseries\color{primary}}{\thesection}{1em}{}
\titleformat{\subsection}{\large\bfseries\color{secondary}}{\thesubsection}{1em}{}
\titleformat{\subsubsection}{\normalsize\bfseries\color{secondary}}{\thesubsubsection}{1em}{}

% Header/Footer
\pagestyle{fancy}
\fancyhf{}
\lhead{\textcolor{primary}{DCCN Lab Viva - Labs 4-7}}
\rhead{\textcolor{primary}{CPE314}}
\rfoot{Page \thepage}

\hypersetup{colorlinks=true,linkcolor=primary,urlcolor=secondary}

% Custom boxes
\newtcolorbox{theorybox}[1]{
    colback=accent!5, colframe=accent,
    title=#1, fonttitle=\bfseries,
    rounded corners, boxrule=1.5pt, breakable
}

\newtcolorbox{commandbox}[1]{
    colback=codebg, colframe=primary,
    title=#1, fonttitle=\bfseries,
    rounded corners, boxrule=1.5pt, breakable
}

\newtcolorbox{notebox}{
    colback=highlight!10, colframe=highlight,
    rounded corners, boxrule=1pt, breakable
}

% Code listing style
\lstset{
    basicstyle=\ttfamily\small,
    backgroundcolor=\color{codebg},
    breaklines=true,
    frame=single,
    rulecolor=\color{primary},
    xleftmargin=10pt,
    xrightmargin=10pt,
    aboveskip=10pt,
    belowskip=10pt,
    columns=flexible
}

\begin{document}

%===============================================================================
% TITLE PAGE
%===============================================================================
\begin{titlepage}
\centering
\vspace*{0.5cm}
{\Large\color{primary}\textbf{DCCN Lab Viva Preparation}}\\[0.3cm]
\rule{\textwidth}{2pt}\\[0.5cm]
{\Huge\bfseries\color{primary} Labs 4, 5, 6, 7}\\[0.3cm]
{\Large Router Configuration \& Routing Protocols}\\[0.2cm]
\rule{\textwidth}{2pt}\\[1.5cm]

\begin{tcolorbox}[colback=light,colframe=primary,rounded corners,width=0.8\textwidth]
\begin{center}
\textbf{\large Lab Topics Covered}\\[0.3cm]
\begin{itemize}[leftmargin=*, itemsep=5pt]
    \item \textbf{Lab 4:} Basic Router Configuration \& Layer 3 Connectivity
    \item \textbf{Lab 5:} Static \& Default Routing Configuration
    \item \textbf{Lab 6:} RIP Dynamic Routing Protocol
    \item \textbf{Lab 7:} OSPF Link-State Routing Protocol
\end{itemize}
\end{center}
\end{tcolorbox}

\vfill
{\large \textbf{Course:} CPE314 - Data Communication and Computer Networks}\\[0.2cm]
{\large \textbf{Tool:} Cisco Packet Tracer}\\[0.3cm]
{\large Spring 2025}
\end{titlepage}

\tableofcontents
\newpage

%===============================================================================
\section{Lab 4: Basic Router Configuration \& Layer 3 Connectivity}
%===============================================================================

\subsection{Objectives}
\begin{itemize}[noitemsep]
    \item To trace configuration procedures for basic layer 3 devices using CLI
    \item To explain and show network connectivity using Cisco Packet Tracer
\end{itemize}

%-------------------------------------------------------------------------------
\subsection{Theory}
%-------------------------------------------------------------------------------

\begin{theorybox}{Cisco Packet Tracer Overview}
\textbf{Cisco Packet Tracer} is a powerful network simulation program that:
\begin{itemize}[noitemsep]
    \item Allows students to experiment with network behavior
    \item Provides simulation, visualization, and assessment capabilities
    \item Enables creation of networks with unlimited devices
    \item Encourages practice, discovery, and troubleshooting
\end{itemize}
\end{theorybox}

\begin{theorybox}{Command Line Interface (CLI) - Operating Modes}
The Cisco IOS has three primary operating modes:

\vspace{0.2cm}
\textbf{1. User EXEC Mode}
\begin{itemize}[noitemsep]
    \item \textbf{Prompt:} \texttt{Router>} or \texttt{Switch>}
    \item \textbf{Access:} Direct login
    \item \textbf{Purpose:} Limited monitoring commands only
    \item \textbf{Limitation:} Cannot configure or restart devices
\end{itemize}

\vspace{0.2cm}
\textbf{2. Privileged EXEC Mode}
\begin{itemize}[noitemsep]
    \item \textbf{Prompt:} \texttt{Router\#} or \texttt{Switch\#}
    \item \textbf{Access:} Enter \texttt{enable} from User EXEC mode
    \item \textbf{Purpose:} View configuration, restart system, enter configuration mode
    \item \textbf{Features:} Access to all router commands
\end{itemize}

\vspace{0.2cm}
\textbf{3. Global Configuration Mode}
\begin{itemize}[noitemsep]
    \item \textbf{Prompt:} \texttt{Router(config)\#}
    \item \textbf{Access:} Enter \texttt{configure terminal} from Privileged mode
    \item \textbf{Purpose:} Commands affect entire system
    \item \textbf{Sub-modes:} Interface config, line config, router config, etc.
\end{itemize}
\end{theorybox}

\begin{theorybox}{Enhanced Editing Commands}
\begin{center}
\begin{tabular}{ll}
\toprule
\textbf{Command} & \textbf{Function} \\
\midrule
\texttt{Ctrl+A} & Move cursor to beginning of line \\
\texttt{Ctrl+E} & Move cursor to end of line \\
\texttt{Esc+B} & Move back one word \\
\texttt{Ctrl+B} & Move back one character \\
\texttt{Ctrl+F} & Move forward one character \\
\texttt{Esc+F} & Move forward one word \\
\texttt{Ctrl+D} / \texttt{Backspace} & Delete a single character \\
\texttt{Ctrl+R} & Redisplay a line \\
\texttt{Ctrl+U} & Erase a line \\
\texttt{Ctrl+W} & Erase a word \\
\texttt{Ctrl+Z} & End configuration mode, return to EXEC \\
\bottomrule
\end{tabular}
\end{center}
\end{theorybox}

\begin{theorybox}{Serial Links: DCE vs DTE}
\textbf{DCE (Data Communication Equipment):}
\begin{itemize}[noitemsep]
    \item Provides clocking signal for synchronization
    \item \textbf{Requires} clock rate configuration
    \item Command: \texttt{clock rate 64000}
\end{itemize}

\vspace{0.2cm}
\textbf{DTE (Data Terminal Equipment):}
\begin{itemize}[noitemsep]
    \item Receives clock signal from DCE
    \item \textbf{Does NOT} require clock rate configuration
\end{itemize}
\end{theorybox}

%-------------------------------------------------------------------------------
\subsection{Configuration Commands by Topic}
%-------------------------------------------------------------------------------

\begin{commandbox}{Process 1: Entering Configuration Mode}
\begin{lstlisting}
Router>enable                           # Enter privileged EXEC mode
Router#configure terminal               # Enter global configuration mode
\end{lstlisting}
\end{commandbox}

\begin{commandbox}{Process 2: Basic Router Configuration}
\begin{lstlisting}
Router(config)#hostname R1              # Set router hostname
R1(config)#no ip domain-lookup          # Disable DNS lookup

# Password Configuration
R1(config)#enable secret class          # Set encrypted EXEC password (preferred)
R1(config)#enable password class        # Set unencrypted EXEC password

# Console Password
R1(config)#line console 0
R1(config-line)#password cisco
R1(config-line)#login
R1(config-line)#exit

# Virtual Terminal (Telnet) Password
R1(config)#line vty 0 4
R1(config-line)#password cisco
R1(config-line)#login
R1(config-line)#exit
\end{lstlisting}
\end{commandbox}

\begin{notebox}
\textbf{Important:} \texttt{enable secret} creates an encrypted password (more secure) while \texttt{enable password} is unencrypted. If both are configured, the secret password takes precedence.
\end{notebox}

\begin{commandbox}{Process 3: FastEthernet Interface Configuration}
\begin{lstlisting}
R1(config)#interface fastethernet 0/0
R1(config-if)#ip address 192.168.1.1 255.255.255.0
R1(config-if)#no shutdown               # Activate interface
R1(config-if)#description Link to LAN   # Optional description
R1(config-if)#exit
\end{lstlisting}
\end{commandbox}

\begin{commandbox}{Process 4: Serial Interface Configuration (DCE Side)}
\begin{lstlisting}
R1(config)#interface serial 0/0/0
R1(config-if)#ip address 192.168.2.1 255.255.255.0
R1(config-if)#clock rate 64000          # Set clock rate (DCE only)
R1(config-if)#no shutdown               # Activate interface
R1(config-if)#description Link to R2    # Optional description
R1(config-if)#exit
\end{lstlisting}
\end{commandbox}

\begin{commandbox}{Process 5: Serial Interface Configuration (DTE Side)}
\begin{lstlisting}
R2(config)#interface serial 0/0/0
R2(config-if)#ip address 192.168.2.2 255.255.255.0
R2(config-if)#no shutdown               # Activate interface
R2(config-if)#exit
# Note: NO clock rate command on DTE side
\end{lstlisting}
\end{commandbox}

\begin{commandbox}{Process 6: Saving Configuration}
\begin{lstlisting}
R1#copy running-config startup-config   # Save to NVRAM (long form)
R1#write memory                         # Save to NVRAM (short form)
R1#wr                                   # Shortest form
\end{lstlisting}
\end{commandbox}

\begin{commandbox}{Process 7: Verification Commands}
\begin{lstlisting}
Router#show ip route                    # Display routing table
Router#show ip interface brief          # Verify interface status
Router#show running-config              # Display current configuration
Router#show startup-config              # Display saved configuration
Router#ping <IP_address>                # Test connectivity (ICMP)
\end{lstlisting}
\end{commandbox}

\begin{commandbox}{Process 8: Configuration Management}
\begin{lstlisting}
Router#erase startup-config             # Delete saved configuration
Router#reload                           # Restart router

# When prompted "Save configuration?", answer "no" to discard changes
\end{lstlisting}
\end{commandbox}

%-------------------------------------------------------------------------------
\subsection{Network Topology Example}
%-------------------------------------------------------------------------------

\begin{center}
\begin{tcolorbox}[colback=light,colframe=primary,width=0.9\textwidth]
\textbf{Typical Lab 4 Topology:}\\[0.2cm]
\texttt{PC1 (192.168.1.10) --- [Fa0/0] R1 [S0/0/0] --- [S0/0/0] R2 [Fa0/0] --- PC2 (192.168.3.10)}\\[0.3cm]
\begin{itemize}[noitemsep]
    \item \textbf{R1 Fa0/0:} 192.168.1.1/24 (LAN side)
    \item \textbf{R1 S0/0/0:} 192.168.2.1/24 (DCE - requires clock rate)
    \item \textbf{R2 S0/0/0:} 192.168.2.2/24 (DTE - no clock rate)
    \item \textbf{R2 Fa0/0:} 192.168.3.1/24 (LAN side)
    \item \textbf{PC1:} 192.168.1.10/24, Gateway 192.168.1.1
    \item \textbf{PC2:} 192.168.3.10/24, Gateway 192.168.3.1
\end{itemize}
\end{tcolorbox}
\end{center}

%-------------------------------------------------------------------------------
\subsection{Key Concepts}
%-------------------------------------------------------------------------------

\begin{itemize}
    \item \textbf{No Shutdown:} Interfaces are administratively down by default; use \texttt{no shutdown} to activate
    \item \textbf{Clock Rate:} Only configured on DCE side of serial links (typically 64000 bps)
    \item \textbf{Administrative Distance:} Directly connected networks have AD = 0 (most trusted)
    \item \textbf{Running-config:} Active configuration in RAM (lost on reboot)
    \item \textbf{Startup-config:} Saved configuration in NVRAM (persistent across reboots)
\end{itemize}

\newpage

%===============================================================================
\section{Lab 5: Static \& Default Routing}
%===============================================================================

\subsection{Objectives}
\begin{itemize}[noitemsep]
    \item Preparation of basic configuration on router based on CLI
    \item Configure and activate serial and Ethernet interfaces
    \item Show connectivity between devices using static routes
    \item Implement default routing for stub networks
\end{itemize}

%-------------------------------------------------------------------------------
\subsection{Theory}
%-------------------------------------------------------------------------------

\begin{theorybox}{Routing Fundamentals}
\textbf{Routing} is the process of taking a packet from one device and sending it through the network to another device on a different network.

\vspace{0.3cm}
\textbf{What Routers Need to Know:}
\begin{itemize}[noitemsep]
    \item Destination address
    \item Neighbor routers from which to learn about remote networks
    \item Possible routes to all remote networks
    \item The best route to each remote network
    \item How to maintain and verify routing information
\end{itemize}
\end{theorybox}

\begin{theorybox}{Types of Routing}
\textbf{1. Static Routing}
\begin{itemize}[noitemsep]
    \item Administrator manually types all network locations into routing table
    \item \textbf{Best for:} Small networks
    \item \textbf{Advantage:} Predictable, secure, no bandwidth overhead
    \item \textbf{Drawback:} Not scalable, requires manual updates
\end{itemize}

\vspace{0.2cm}
\textbf{2. Dynamic Routing}
\begin{itemize}[noitemsep]
    \item Protocols automatically update routing tables
    \item \textbf{Best for:} Larger networks
    \item \textbf{Advantage:} Automatic updates, adapts to topology changes
    \item \textbf{Drawback:} Requires CPU, memory, and bandwidth overhead
\end{itemize}
\end{theorybox}

\begin{theorybox}{Administrative Distances}
\textbf{AD} is a measure of trustworthiness (lower = more trusted):

\begin{center}
\begin{tabular}{ll}
\toprule
\textbf{Route Source} & \textbf{Default AD} \\
\midrule
Connected interface & 0 (most trusted) \\
Static route & 1 \\
EIGRP & 90 \\
OSPF & 110 \\
RIP & 120 (least trusted) \\
\bottomrule
\end{tabular}
\end{center}

\textbf{Metric for Static Routes:} 0 (administrator decision, not protocol)
\end{theorybox}

\begin{theorybox}{Default Route Concept}
A \textbf{default route} (0.0.0.0/0) is a catch-all route used when no specific route exists for a destination.

\vspace{0.2cm}
\textbf{Use Case:} Ideal for stub routers (routers with only one exit path)

\vspace{0.2cm}
\textbf{Benefit:} Reduces routing table size
\end{theorybox}

%-------------------------------------------------------------------------------
\subsection{Configuration Commands by Topic}
%-------------------------------------------------------------------------------

\begin{commandbox}{Process 1: Static Route Configuration}
\textbf{Syntax:}
\begin{lstlisting}
Router(config)#ip route <network-address> <subnet-mask> <next-hop-ip>
\end{lstlisting}

\textbf{Where:}
\begin{itemize}[noitemsep]
    \item \texttt{<network-address>} = Destination network address
    \item \texttt{<subnet-mask>} = Subnet mask of remote network
    \item \texttt{<next-hop-ip>} = IP address of next-hop router
\end{itemize}

\vspace{0.2cm}
\textbf{Example: R1 needs route to Network 192.168.3.0/24 via R2 (192.168.2.2)}
\begin{lstlisting}
R1(config)#ip route 192.168.3.0 255.255.255.0 192.168.2.2
\end{lstlisting}
\end{commandbox}

\begin{commandbox}{Process 2: Static Route Configuration (Multiple Networks)}
\textbf{Scenario:} Three routers (R1 --- R2 --- R3) with three networks

\vspace{0.2cm}
\textbf{On R1 (leftmost):}
\begin{lstlisting}
R1(config)#ip route 192.168.2.0 255.255.255.0 <R2_next_hop>
R1(config)#ip route 192.168.3.0 255.255.255.0 <R2_next_hop>
\end{lstlisting}

\textbf{On R2 (center):}
\begin{lstlisting}
R2(config)#ip route 192.168.1.0 255.255.255.0 <R1_next_hop>
R2(config)#ip route 192.168.3.0 255.255.255.0 <R3_next_hop>
\end{lstlisting}

\textbf{On R3 (rightmost):}
\begin{lstlisting}
R3(config)#ip route 192.168.1.0 255.255.255.0 <R2_next_hop>
R3(config)#ip route 192.168.2.0 255.255.255.0 <R2_next_hop>
\end{lstlisting}
\end{commandbox}

\begin{commandbox}{Process 3: Default Static Route Configuration}
\textbf{Syntax:}
\begin{lstlisting}
Router(config)#ip route 0.0.0.0 0.0.0.0 <next-hop-ip>
\end{lstlisting}

\textbf{Explanation:}
\begin{itemize}[noitemsep]
    \item \texttt{0.0.0.0} = Match any destination network
    \item \texttt{0.0.0.0} = Match any subnet mask
    \item Used for stub routers with single exit point
\end{itemize}

\vspace{0.2cm}
\textbf{Example: R1 uses R2 as default gateway}
\begin{lstlisting}
R1(config)#ip route 0.0.0.0 0.0.0.0 192.168.2.2
\end{lstlisting}
\end{commandbox}

\begin{commandbox}{Process 4: Removing Static Routes}
\begin{lstlisting}
Router(config)#no ip route <network> <subnet-mask> <next-hop>

# Example: Remove specific static route on R1
R1(config)#no ip route 192.168.3.0 255.255.255.0 192.168.2.2
\end{lstlisting}
\end{commandbox}

\begin{commandbox}{Process 5: Verification Commands}
\begin{lstlisting}
Router#show ip route                    # Display routing table
Router#show ip interface brief          # Verify interface status
Router#ping <IP_address>                # Test connectivity

# In routing table output:
#   S = Static route
#   C = Connected network
#   * = Default route candidate
\end{lstlisting}
\end{commandbox}

%-------------------------------------------------------------------------------
\subsection{Network Topology Example}
%-------------------------------------------------------------------------------

\begin{center}
\begin{tcolorbox}[colback=light,colframe=primary,width=0.9\textwidth]
\textbf{Typical Lab 5 Topology (Three Routers):}\\[0.2cm]
\texttt{PC1 --- R1 --- R2 --- R3 --- PC3}\\[0.3cm]
\begin{itemize}[noitemsep]
    \item \textbf{Network 1 (R1 LAN):} 192.168.1.0/24
    \item \textbf{Network 2 (R1-R2 link):} 192.168.2.0/24
    \item \textbf{Network 3 (R2-R3 link):} 192.168.3.0/24
    \item \textbf{Network 4 (R3 LAN):} 192.168.4.0/24
\end{itemize}

\vspace{0.2cm}
\textbf{Configuration Strategy:}
\begin{itemize}[noitemsep]
    \item R1 and R3 are stub routers $\rightarrow$ Use default routes
    \item R2 is the core router $\rightarrow$ Needs specific static routes to all networks
\end{itemize}
\end{tcolorbox}
\end{center}

%-------------------------------------------------------------------------------
\subsection{Key Concepts}
%-------------------------------------------------------------------------------

\begin{itemize}
    \item \textbf{Stub Router:} Router with only one exit path; ideal for default routes
    \item \textbf{Next-Hop Address:} IP address of neighboring router's interface
    \item \textbf{Static Route Metric:} Always 0 (manually configured, no protocol metric)
    \item \textbf{Route Prefix 'S':} Indicates static route in routing table
    \item \textbf{Manual Configuration:} Static routes must be configured on each router separately
\end{itemize}

\newpage

%===============================================================================
\section{Lab 6: RIP Dynamic Routing Protocol}
%===============================================================================

\subsection{Objectives}
\begin{itemize}[noitemsep]
    \item To construct a network demonstrating RIP routing protocol operation
    \item To show connectivity between nodes using RIP on all routers
    \item Understand distance-vector routing protocol behavior
\end{itemize}

%-------------------------------------------------------------------------------
\subsection{Theory}
%-------------------------------------------------------------------------------

\begin{theorybox}{Dynamic Routing Protocols Overview}
\textbf{Dynamic Routing} is when protocols automatically discover networks and update routing tables on routers. Routers communicate routing information with neighboring routers.

\vspace{0.3cm}
\textbf{Advantages:}
\begin{itemize}[noitemsep]
    \item Automatic network discovery
    \item Adapts to topology changes
    \item Scalable for larger networks
\end{itemize}
\end{theorybox}

\begin{theorybox}{Two Categories of Dynamic Routing}
\textbf{1. Distance-Vector Routing}
\begin{itemize}[noitemsep]
    \item Finds best path by judging distance (hop count)
    \item Periodically sends entire routing table to neighbors
    \item \textbf{Example:} RIP (Routing Information Protocol)
\end{itemize}

\vspace{0.2cm}
\textbf{2. Link-State Routing}
\begin{itemize}[noitemsep]
    \item Creates three tables: neighbors, topology, routing
    \item Sends updates containing state of own links to all routers
    \item \textbf{Example:} OSPF (Open Shortest Path First)
\end{itemize}
\end{theorybox}

\begin{theorybox}{Routing Information Protocol (RIP) Characteristics}
\begin{itemize}[noitemsep]
    \item \textbf{Type:} Distance-vector routing protocol
    \item \textbf{Metric:} Hop count (number of routers between source and destination)
    \item \textbf{Maximum Hop Count:} 15 (16 = unreachable)
    \item \textbf{Update Frequency:} Sends complete routing table every 30 seconds
    \item \textbf{Administrative Distance:} 120 (least trusted)
    \item \textbf{Use Case:} Small networks only
    \item \textbf{Limitations:} Inefficient on large networks with slow WAN links
\end{itemize}
\end{theorybox}

\begin{theorybox}{RIP Version Comparison}
\begin{center}
\begin{tabular}{lll}
\toprule
\textbf{Feature} & \textbf{RIPv1} & \textbf{RIPv2} \\
\midrule
Routing Type & Classful & Classless \\
Subnet Mask Info & Not sent & Sent in updates \\
VLSM Support & No & Yes \\
Modern Use & Rarely used & More common \\
\bottomrule
\end{tabular}
\end{center}
\end{theorybox}

%-------------------------------------------------------------------------------
\subsection{Configuration Commands by Topic}
%-------------------------------------------------------------------------------

\begin{commandbox}{Process 1: Enabling RIP on Router}
\textbf{Syntax:}
\begin{lstlisting}
Router(config)#router rip
Router(config-router)#network <NETWORK_ADDRESS>
Router(config-router)#network <NETWORK_ADDRESS>
Router(config-router)#exit
\end{lstlisting}

\textbf{Important Notes:}
\begin{itemize}[noitemsep]
    \item Enter \textbf{each directly connected network} separately
    \item Use \textbf{classful network addresses} (not subnet addresses)
    \item RIP automatically determines which interfaces belong to each network
\end{itemize}
\end{commandbox}

\begin{commandbox}{Process 2: RIP Configuration Example (Router R1)}
\textbf{Scenario:} R1 has two interfaces:
\begin{itemize}[noitemsep]
    \item Fa0/0: 10.1.1.1/24 (connected to Network 10.0.0.0)
    \item S0/0/0: 10.2.2.1/24 (connected to Network 10.0.0.0)
\end{itemize}

\begin{lstlisting}
R1(config)#router rip
R1(config-router)#network 10.0.0.0      # Classful network address
R1(config-router)#exit
\end{lstlisting}

\textbf{Note:} Since both interfaces are in Class A network 10.0.0.0, only one \texttt{network} command needed.
\end{commandbox}

\begin{commandbox}{Process 3: RIP Configuration (Multiple Networks)}
\textbf{Scenario:} Router R2 connects three networks:
\begin{itemize}[noitemsep]
    \item Fa0/0: 172.16.1.1/24
    \item S0/0/0: 172.17.1.1/24
    \item S0/0/1: 172.18.1.1/24
\end{itemize}

\begin{lstlisting}
R2(config)#router rip
R2(config-router)#network 172.16.0.0   # Network 1
R2(config-router)#network 172.17.0.0   # Network 2
R2(config-router)#network 172.18.0.0   # Network 3
R2(config-router)#exit
\end{lstlisting}
\end{commandbox}

\begin{commandbox}{Process 4: Verifying RIP Configuration}
\begin{lstlisting}
Router#show ip route                    # Display routing table
Router#show ip protocols                # Display routing protocol info
Router#show running-config              # Verify RIP configuration
Router#debug ip rip                     # Show RIP updates in real-time
Router#no debug ip rip                  # Turn off RIP debugging
\end{lstlisting}

\textbf{Interpreting Routing Table:}
\begin{itemize}[noitemsep]
    \item \textbf{R} = Route learned via RIP
    \item \textbf{[120/X]} = Administrative Distance 120, Metric X hops
    \item \textbf{via <IP>} = Next-hop router address
\end{itemize}
\end{commandbox}

\begin{commandbox}{Process 5: Disabling RIP}
\begin{lstlisting}
Router(config)#no router rip            # Completely remove RIP

# Remove specific network from RIP
Router(config)#router rip
Router(config-router)#no network <NETWORK_ADDRESS>
\end{lstlisting}
\end{commandbox}

%-------------------------------------------------------------------------------
\subsection{Network Topology Example}
%-------------------------------------------------------------------------------

\begin{center}
\begin{tcolorbox}[colback=light,colframe=primary,width=0.9\textwidth]
\textbf{Typical Lab 6 Topology (Class A Addressing):}\\[0.2cm]
\texttt{PC1 --- R1 --- R2 --- R3 --- PC3}\\[0.3cm]
\begin{itemize}[noitemsep]
    \item \textbf{Network 1 (R1 LAN):} 10.1.0.0/24
    \item \textbf{Network 2 (R1-R2 link):} 10.2.0.0/24
    \item \textbf{Network 3 (R2 LAN):} 10.3.0.0/24
    \item \textbf{Network 4 (R2-R3 link):} 10.4.0.0/24
    \item \textbf{Network 5 (R3 LAN):} 10.5.0.0/24
\end{itemize}

\vspace{0.2cm}
\textbf{Configuration Sequence:}
\begin{enumerate}[noitemsep]
    \item Enable RIP on R1 (networks 10.0.0.0)
    \item Enable RIP on R3 (networks 10.0.0.0)
    \item Enable RIP on R2 (networks 10.0.0.0) $\rightarrow$ Full connectivity achieved
\end{enumerate}
\end{tcolorbox}
\end{center}

%-------------------------------------------------------------------------------
\subsection{Key Concepts}
%-------------------------------------------------------------------------------

\begin{itemize}
    \item \textbf{Convergence:} Time taken for all routers to learn all routes (30+ seconds for RIP)
    \item \textbf{Hop Count:} Metric used by RIP; each router hop = 1
    \item \textbf{15-Hop Limit:} Maximum reachable distance; 16 hops = unreachable
    \item \textbf{Periodic Updates:} RIP sends updates every 30 seconds
    \item \textbf{Administrative Distance 120:} RIP is least trusted routing protocol
    \item \textbf{Classful Addressing:} RIPv1 uses classful networks (Class A, B, C)
    \item \textbf{Network Discovery:} RIP automatically learns remote networks from neighbors
\end{itemize}

\newpage

%===============================================================================
\section{Lab 7: OSPF Link-State Routing Protocol}
%===============================================================================

\subsection{Objectives}
\begin{itemize}[noitemsep]
    \item To construct a network demonstrating OSPF routing protocol operation
    \item To show connectivity between nodes using OSPF on all routers
    \item Understand link-state routing protocol behavior
    \item Configure OSPF areas, router IDs, and interface parameters
\end{itemize}

%-------------------------------------------------------------------------------
\subsection{Theory}
%-------------------------------------------------------------------------------

\begin{theorybox}{Open Shortest Path First (OSPF) Characteristics}
\begin{itemize}[noitemsep]
    \item \textbf{Type:} Link-state routing protocol
    \item \textbf{Standard:} Open standard (multi-vendor support)
    \item \textbf{Algorithm:} Dijkstra shortest path first (SPF)
    \item \textbf{Metric:} Cost (based on bandwidth)
    \item \textbf{Administrative Distance:} 110
    \item \textbf{Convergence:} Fast convergence to topology changes
    \item \textbf{Scalability:} Supports unlimited hop count
    \item \textbf{Features:} Supports VLSM, CIDR, hierarchical design with areas
\end{itemize}
\end{theorybox}

\begin{theorybox}{OSPF Advantages Over RIP}
\begin{itemize}[noitemsep]
    \item \textbf{Unlimited Hop Count:} No 15-hop limitation
    \item \textbf{Fast Convergence:} Immediately detects topology changes
    \item \textbf{Efficient Updates:} Sends only changes, not entire routing table
    \item \textbf{Hierarchical Design:} Supports areas for better scalability
    \item \textbf{Load Balancing:} Equal-cost multipath routing
    \item \textbf{VLSM Support:} Variable length subnet masking
\end{itemize}
\end{theorybox}

\begin{theorybox}{OSPF Metric: Cost Calculation}
\textbf{Cost Formula:}
\begin{center}
\large{$\text{Cost} = \frac{10^8}{\text{Bandwidth (bps)}}$}
\end{center}

\vspace{0.2cm}
\textbf{Examples:}
\begin{center}
\begin{tabular}{lll}
\toprule
\textbf{Interface Type} & \textbf{Bandwidth} & \textbf{Cost} \\
\midrule
10 Mbps Ethernet & 10,000,000 bps & 10 \\
100 Mbps FastEthernet & 100,000,000 bps & 1 \\
1 Gbps GigabitEthernet & 1,000,000,000 bps & 1 \\
64 kbps Serial & 64,000 bps & 1562 \\
\bottomrule
\end{tabular}
\end{center}

\textbf{Best Path:} OSPF selects route with \textbf{lowest cumulative cost}
\end{theorybox}

\begin{theorybox}{Router ID (RID)}
\textbf{Definition:} IP address used to uniquely identify an OSPF router

\vspace{0.2cm}
\textbf{Selection Precedence (highest to lowest):}
\begin{enumerate}[noitemsep]
    \item IP address configured with \texttt{router-id} command
    \item Highest IP address of active \textbf{loopback} interfaces
    \item Highest IP address of active \textbf{physical} interfaces
\end{enumerate}

\vspace{0.2cm}
\textbf{Best Practice:} Configure loopback interface for stable Router ID
\end{theorybox}

\begin{theorybox}{OSPF Terminology}
\textbf{Neighbor:}
\begin{itemize}[noitemsep]
    \item Two or more routers with interfaces on common network
    \item Discovered via Hello protocol
\end{itemize}

\vspace{0.2cm}
\textbf{Adjacency:}
\begin{itemize}[noitemsep]
    \item Relationship between routers permitting direct route exchange
    \item Not all neighbors become adjacent
\end{itemize}

\vspace{0.2cm}
\textbf{Designated Router (DR):}
\begin{itemize}[noitemsep]
    \item Elected on broadcast/multi-access networks
    \item Minimizes number of adjacencies
    \item Election based on: Highest priority $\rightarrow$ Highest Router ID
\end{itemize}

\vspace{0.2cm}
\textbf{Backup Designated Router (BDR):}
\begin{itemize}[noitemsep]
    \item Hot standby for DR
    \item Receives all routing updates but doesn't disperse until DR fails
\end{itemize}

\vspace{0.2cm}
\textbf{Hello Protocol:}
\begin{itemize}[noitemsep]
    \item Provides dynamic neighbor discovery
    \item Maintains neighbor relationships
    \item Multicast address: 224.0.0.5
    \item Default intervals: Hello = 10 sec, Dead = 40 sec
\end{itemize}

\vspace{0.2cm}
\textbf{Link State Advertisement (LSA):}
\begin{itemize}[noitemsep]
    \item OSPF data packet containing link-state information
    \item Shared between adjacent routers
\end{itemize}

\vspace{0.2cm}
\textbf{Topological Database:}
\begin{itemize}[noitemsep]
    \item Contains information from all LSA packets
    \item Input for Dijkstra SPF algorithm
\end{itemize}
\end{theorybox}

\begin{theorybox}{OSPF Areas}
\textbf{Definition:} Grouping of contiguous networks and routers

\vspace{0.2cm}
\textbf{Purpose:}
\begin{itemize}[noitemsep]
    \item Reduce routing overhead
    \item Speed up convergence
    \item Confine network instability to single areas
\end{itemize}

\vspace{0.2cm}
\textbf{Area 0 (Backbone Area):}
\begin{itemize}[noitemsep]
    \item Required in multi-area OSPF
    \item All other areas must connect to Area 0
\end{itemize}

\vspace{0.2cm}
\textbf{Single Area:} For simple topologies, use Area 0 for entire network
\end{theorybox}

\begin{theorybox}{OSPF Network Types}
\begin{enumerate}[noitemsep]
    \item \textbf{Broadcast (Multi-access):} Ethernet - DR/BDR elected
    \item \textbf{Non-Broadcast Multi-access (NBMA):} Frame Relay - special config
    \item \textbf{Point-to-Point:} Serial links - no DR/BDR needed
    \item \textbf{Point-to-Multipoint:} Hub-and-spoke topologies
\end{enumerate}
\end{theorybox}

%-------------------------------------------------------------------------------
\subsection{Configuration Commands by Topic}
%-------------------------------------------------------------------------------

\begin{commandbox}{Process 1: Basic OSPF Configuration}
\textbf{Syntax:}
\begin{lstlisting}
Router(config)#router ospf <Process_ID>
Router(config-router)#network <Network_Address> <Wildcard_Mask> area <Area_ID>
Router(config-router)#network <Network_Address> <Wildcard_Mask> area <Area_ID>
Router(config-router)#exit
\end{lstlisting}

\textbf{Parameters:}
\begin{itemize}[noitemsep]
    \item \texttt{<Process\_ID>} = Locally significant number (1-65535)
    \item \texttt{<Network\_Address>} = Network or IP address to advertise
    \item \texttt{<Wildcard\_Mask>} = Inverse of subnet mask
    \item \texttt{<Area\_ID>} = OSPF area number (typically 0)
\end{itemize}
\end{commandbox}

\begin{notebox}
\textbf{Wildcard Mask Calculation:}\\
Wildcard Mask = 255.255.255.255 - Subnet Mask

\vspace{0.2cm}
\textbf{Examples:}
\begin{itemize}[noitemsep]
    \item Subnet Mask 255.255.255.0 $\rightarrow$ Wildcard 0.0.0.255
    \item Subnet Mask 255.255.255.252 $\rightarrow$ Wildcard 0.0.0.3
    \item Subnet Mask 255.255.255.128 $\rightarrow$ Wildcard 0.0.0.127
\end{itemize}
\end{notebox}

\begin{commandbox}{Process 2: OSPF Configuration Example (Router R1)}
\textbf{Scenario:} R1 has three interfaces in Class B network (172.17.0.0):
\begin{itemize}[noitemsep]
    \item Fa0/0: 172.17.1.1/24
    \item S0/0/0: 172.17.2.1/30
    \item S0/0/1: 172.17.3.1/30
\end{itemize}

\begin{lstlisting}
R1(config)#router ospf 1
R1(config-router)#network 172.17.1.0 0.0.0.255 area 0
R1(config-router)#network 172.17.2.0 0.0.0.3 area 0
R1(config-router)#network 172.17.3.0 0.0.0.3 area 0
R1(config-router)#exit
\end{lstlisting}
\end{commandbox}

\begin{commandbox}{Process 3: Configuring Router ID with Loopback}
\textbf{Purpose:} Provide stable, predictable Router ID

\begin{lstlisting}
R1(config)#interface loopback 0
R1(config-if)#ip address 10.1.1.1 255.255.255.255
R1(config-if)#exit

R2(config)#interface loopback 0
R2(config-if)#ip address 10.2.2.2 255.255.255.255
R2(config-if)#exit

R3(config)#interface loopback 0
R3(config-if)#ip address 10.3.3.3 255.255.255.255
R3(config-if)#exit

# Reload router for new Router ID to take effect
Router#reload
\end{lstlisting}

\textbf{Note:} Use /32 mask (255.255.255.255) for loopback interfaces
\end{commandbox}

\begin{commandbox}{Process 4: Configuring Interface Bandwidth}
\textbf{Purpose:} Accurate OSPF cost calculation

\begin{lstlisting}
R1(config)#interface serial 0/0/0
R1(config-if)#bandwidth 64         # Set bandwidth to 64 kbps
R1(config-if)#exit

R2(config)#interface serial 0/0/0
R2(config-if)#bandwidth 128        # Set bandwidth to 128 kbps
R2(config-if)#exit
\end{lstlisting}

\textbf{Note:} This command does NOT change actual interface speed; only affects metric calculation
\end{commandbox}

\begin{commandbox}{Process 5: Modifying OSPF Hello and Dead Intervals}
\textbf{Default Values:}
\begin{itemize}[noitemsep]
    \item Hello Interval: 10 seconds
    \item Dead Interval: 40 seconds (4 $\times$ Hello)
\end{itemize}

\begin{lstlisting}
R1(config)#interface serial 0/0/0
R1(config-if)#ip ospf hello-interval 5     # Hello every 5 seconds
R1(config-if)#ip ospf dead-interval 20     # Dead after 20 seconds
R1(config-if)#exit

# MUST match on both sides of link
R2(config)#interface serial 0/0/0
R2(config-if)#ip ospf hello-interval 5
R2(config-if)#ip ospf dead-interval 20
R2(config-if)#exit
\end{lstlisting}

\textbf{Important:} Intervals must match on both routers; mismatch prevents adjacency
\end{commandbox}

\begin{commandbox}{Process 6: Configuring OSPF Priority (DR/BDR Election)}
\begin{lstlisting}
R1(config)#interface fastethernet 0/0
R1(config-if)#ip ospf priority 255    # Highest priority = DR
R1(config-if)#exit

R2(config)#interface fastethernet 0/0
R2(config-if)#ip ospf priority 100    # Medium priority = BDR
R2(config-if)#exit

R3(config)#interface fastethernet 0/0
R3(config-if)#ip ospf priority 0      # Priority 0 = Never DR/BDR
R3(config-if)#exit
\end{lstlisting}

\textbf{Default Priority:} 1\\
\textbf{Range:} 0-255 (higher = better chance to become DR)
\end{commandbox}

\begin{commandbox}{Process 7: OSPF Verification Commands}
\begin{lstlisting}
Router#show ip protocols                # Display routing protocol info
Router#show ip ospf                     # Display OSPF process details
Router#show ip ospf neighbor            # Display OSPF neighbors
Router#show ip ospf interface           # Display OSPF interface details
Router#show ip ospf interface brief     # Brief OSPF interface summary
Router#show ip route                    # Display routing table
Router#show ip route ospf               # Display only OSPF routes
Router#debug ip ospf events             # Show OSPF events
Router#debug ip ospf hello              # Show Hello packets
Router#no debug all                     # Turn off all debugging
\end{lstlisting}

\textbf{Interpreting Routing Table:}
\begin{itemize}[noitemsep]
    \item \textbf{O} = OSPF intra-area route
    \item \textbf{[110/X]} = AD 110, Cost X
    \item \textbf{via <IP>} = Next-hop router address
\end{itemize}
\end{commandbox}

%-------------------------------------------------------------------------------
\subsection{Network Topology Example}
%-------------------------------------------------------------------------------

\begin{center}
\begin{tcolorbox}[colback=light,colframe=primary,width=0.95\textwidth]
\textbf{Typical Lab 7 Topology (Class B with VLSM):}\\[0.2cm]
\texttt{PC1 --- R1 --- R2 --- R3 --- PC3}\\[0.3cm]
\begin{itemize}[noitemsep]
    \item \textbf{Network 1 (R1 LAN):} 172.17.1.0/24
    \item \textbf{Network 2 (R1-R2 link):} 172.17.2.0/30 (point-to-point)
    \item \textbf{Network 3 (R2 LAN):} 172.17.3.0/24
    \item \textbf{Network 4 (R2-R3 link):} 172.17.4.0/30 (point-to-point)
    \item \textbf{Network 5 (R3 LAN):} 172.17.5.0/24
    \item \textbf{Loopbacks:} R1=10.1.1.1, R2=10.2.2.2, R3=10.3.3.3
\end{itemize}

\vspace{0.2cm}
\textbf{Configuration Order:}
\begin{enumerate}[noitemsep]
    \item Configure loopback interfaces on all routers
    \item Enable OSPF process 1 on all routers
    \item Add all directly connected networks to Area 0
    \item Verify neighbor adjacencies
    \item Test connectivity between all hosts
\end{enumerate}
\end{tcolorbox}
\end{center}

%-------------------------------------------------------------------------------
\subsection{Key Concepts}
%-------------------------------------------------------------------------------

\begin{itemize}
    \item \textbf{Process ID:} Locally significant; can differ between routers
    \item \textbf{Area ID:} Must match for routers to become neighbors
    \item \textbf{Wildcard Mask:} Inverse of subnet mask (used for network matching)
    \item \textbf{Cost Metric:} Based on bandwidth; lower cost = better path
    \item \textbf{Router ID:} Loopback interface provides stable RID
    \item \textbf{Administrative Distance:} 110 (more trusted than RIP's 120)
    \item \textbf{Convergence:} Faster than RIP; immediate topology change detection
    \item \textbf{Hello Timers:} Must match on both sides for adjacency
    \item \textbf{DR/BDR:} Only on broadcast networks; not needed on point-to-point links
    \item \textbf{Area 0:} Backbone area; single-area designs use Area 0 for all networks
\end{itemize}

\newpage

%===============================================================================
\section{Comparison Summary}
%===============================================================================

\begin{center}
\begin{tcolorbox}[colback=light,colframe=primary,title=\textbf{\large Routing Methods Comparison},width=0.95\textwidth]
\begin{center}
\begin{tabular}{p{3cm}p{3.5cm}p{3.5cm}p{3.5cm}}
\toprule
\textbf{Feature} & \textbf{Static Routing} & \textbf{RIP} & \textbf{OSPF} \\
\midrule
\textbf{Type} & Manual configuration & Distance-vector & Link-state \\[0.2cm]
\textbf{Metric} & N/A (admin decision) & Hop count & Cost (bandwidth) \\[0.2cm]
\textbf{Admin Distance} & 1 & 120 & 110 \\[0.2cm]
\textbf{Convergence} & N/A & Slow (30+ sec) & Fast (immediate) \\[0.2cm]
\textbf{Scalability} & Small networks & Small networks & Large networks \\[0.2cm]
\textbf{Configuration} & Manual per route & Network statements & Network + area \\[0.2cm]
\textbf{Updates} & None & Full table every 30s & Only changes \\[0.2cm]
\textbf{Hop Limit} & Unlimited & 15 hops max & Unlimited \\[0.2cm]
\textbf{VLSM Support} & Yes & No (v1) / Yes (v2) & Yes \\[0.2cm]
\textbf{Bandwidth Usage} & None & High & Low \\[0.2cm]
\textbf{Best Use} & Stub networks & Legacy small nets & Modern enterprise \\[0.2cm]
\bottomrule
\end{tabular}
\end{center}
\end{tcolorbox}
\end{center}

\vspace{0.5cm}

\begin{center}
\begin{tcolorbox}[colback=accent!10,colframe=accent,title=\textbf{\large Command Syntax Quick Reference},width=0.95\textwidth]
\textbf{Static Route:}
\begin{lstlisting}
ip route <destination-network> <mask> <next-hop-ip>
\end{lstlisting}

\textbf{Default Route:}
\begin{lstlisting}
ip route 0.0.0.0 0.0.0.0 <next-hop-ip>
\end{lstlisting}

\textbf{RIP:}
\begin{lstlisting}
router rip
network <classful-network-address>
\end{lstlisting}

\textbf{OSPF:}
\begin{lstlisting}
router ospf <process-id>
network <network-address> <wildcard-mask> area <area-id>
\end{lstlisting}
\end{tcolorbox}
\end{center}

%===============================================================================
\section{Common Verification Commands}
%===============================================================================

\begin{center}
\begin{tcolorbox}[colback=light,colframe=secondary,width=0.95\textwidth]
\textbf{Universal Commands (All Labs):}
\begin{lstlisting}
show ip route                           # Routing table
show ip interface brief                 # Interface status summary
show running-config                     # Current configuration
ping <IP>                               # Test connectivity
traceroute <IP>                         # Trace packet path
\end{lstlisting}

\vspace{0.3cm}
\textbf{RIP-Specific:}
\begin{lstlisting}
show ip protocols                       # Routing protocol details
debug ip rip                            # Show RIP updates
\end{lstlisting}

\vspace{0.3cm}
\textbf{OSPF-Specific:}
\begin{lstlisting}
show ip ospf neighbor                   # OSPF neighbors
show ip ospf interface                  # OSPF interface details
show ip ospf                            # OSPF process information
show ip route ospf                      # Only OSPF routes
debug ip ospf events                    # OSPF events
debug ip ospf hello                     # Hello packets
\end{lstlisting}
\end{tcolorbox}
\end{center}

%===============================================================================
\section{Troubleshooting Tips}
%===============================================================================

\subsection{Lab 4: Basic Configuration Issues}
\begin{itemize}
    \item \textbf{Interface down:} Check \texttt{no shutdown} command
    \item \textbf{No connectivity:} Verify IP addressing and subnet masks
    \item \textbf{Serial link down:} Verify clock rate on DCE side
    \item \textbf{Cannot ping:} Check default gateways on PCs
\end{itemize}

\subsection{Lab 5: Static Routing Issues}
\begin{itemize}
    \item \textbf{No route in table:} Check static route syntax
    \item \textbf{Wrong next-hop:} Verify next-hop IP is reachable
    \item \textbf{Default route not working:} Ensure 0.0.0.0/0 configured correctly
\end{itemize}

\subsection{Lab 6: RIP Issues}
\begin{itemize}
    \item \textbf{No RIP routes:} Verify \texttt{router rip} and \texttt{network} commands
    \item \textbf{Slow convergence:} Wait 30+ seconds for updates
    \item \textbf{Missing networks:} Check classful network addresses used
    \item \textbf{Hop count exceeded:} Destination $>$15 hops away
\end{itemize}

\subsection{Lab 7: OSPF Issues}
\begin{itemize}
    \item \textbf{No neighbors:} Check area ID, Hello/Dead timers match
    \item \textbf{Wrong Router ID:} Configure loopback and reload
    \item \textbf{High cost routes:} Verify bandwidth settings
    \item \textbf{No adjacency:} Check wildcard masks, area numbers
    \item \textbf{Stuck in INIT state:} Check bidirectional communication
\end{itemize}

%===============================================================================
\section{Viva Preparation Checklist}
%===============================================================================

\begin{tcolorbox}[colback=highlight!15,colframe=highlight,title=\textbf{Before Your Viva},rounded corners,breakable]
\textbf{Conceptual Understanding:}
\begin{itemize}[noitemsep]
    \item[$\Box$] Explain difference between static and dynamic routing
    \item[$\Box$] Describe how distance-vector routing works
    \item[$\Box$] Describe how link-state routing works
    \item[$\Box$] Compare RIP vs OSPF characteristics
    \item[$\Box$] Explain administrative distance concept
    \item[$\Box$] Understand DCE vs DTE on serial links
\end{itemize}

\vspace{0.3cm}
\textbf{Practical Skills:}
\begin{itemize}[noitemsep]
    \item[$\Box$] Configure router hostname, passwords
    \item[$\Box$] Configure FastEthernet and Serial interfaces
    \item[$\Box$] Create static routes with correct syntax
    \item[$\Box$] Configure default routes for stub routers
    \item[$\Box$] Enable RIP with network statements
    \item[$\Box$] Configure OSPF with wildcard masks
    \item[$\Box$] Set up loopback interfaces for Router ID
\end{itemize}

\vspace{0.3cm}
\textbf{Verification \& Troubleshooting:}
\begin{itemize}[noitemsep]
    \item[$\Box$] Use \texttt{show ip route} effectively
    \item[$\Box$] Interpret routing table entries
    \item[$\Box$] Verify neighbor adjacencies (OSPF)
    \item[$\Box$] Test connectivity with ping
    \item[$\Box$] Calculate wildcard masks
    \item[$\Box$] Calculate OSPF cost
\end{itemize}
\end{tcolorbox}

\vfill

\begin{center}
\begin{tcolorbox}[colback=primary!10,colframe=primary,width=0.7\textwidth]
\begin{center}
\textbf{\large Good Luck with Your Viva!}\\[0.2cm]
\textit{Practice makes perfect}
\end{center}
\end{tcolorbox}
\end{center}

\end{document}
